% ============================================================
%  Topological Selection of the Inflationary Winding Number and
%  CMB Predictions from Braided Kaluza-Klein Geometry
%
%  arXiv submission — primary: hep-th   cross-list: gr-qc, astro-ph.CO
%  Target: ~20 pages, self-contained
% ============================================================
\documentclass[12pt,a4paper]{article}

% --- Packages ---
\usepackage{amsmath,amssymb,amsthm}
\usepackage{geometry}
\geometry{margin=2.5cm}
\usepackage{hyperref}
\hypersetup{
  colorlinks=true,
  linkcolor=blue,
  citecolor=blue,
  urlcolor=blue,
  pdftitle={Topological Selection of the Inflationary Winding Number and CMB Predictions from Braided Kaluza-Klein Geometry},
  pdfauthor={ThomasCory Walker-Pearson}
}
\usepackage{booktabs}
\usepackage{mathtools}
\usepackage{natbib}
\bibliographystyle{unsrtnat}

% --- Theorem environments ---
\newtheorem{theorem}{Theorem}[section]
\newtheorem{proposition}[theorem]{Proposition}
\newtheorem{lemma}[theorem]{Lemma}
\newtheorem{corollary}[theorem]{Corollary}
\newtheorem{remark}[theorem]{Remark}

% --- Macros ---
\newcommand{\nw}{n_{\mathrm{w}}}
\newcommand{\kcs}{k_{\mathrm{CS}}}
\newcommand{\cs}{c_{\mathrm{s}}}
\newcommand{\phiz}{\phi_0}
\newcommand{\ns}{n_{\mathrm{s}}}
\newcommand{\rten}{r}
\newcommand{\Bmu}{B_{\mu}}
\newcommand{\lam}{\lambda}
\newcommand{\aNM}{\alpha_{\mathrm{NM}}}
\newcommand{\aem}{\alpha_{\mathrm{em}}}
\newcommand{\lP}{\ell_P}
\newcommand{\Gfive}{G_{AB}}

% ============================================================
\begin{document}

\title{\textbf{Topological Selection of the Inflationary Winding Number}\\
\textbf{and CMB Predictions from Braided Kaluza-Klein Geometry}}

\author{ThomasCory Walker-Pearson\\[4pt]
\small AxiomZero Technologies (DBA) / Independent Researcher, Pacific Northwest, USA\\[2pt]
\small \texttt{https://github.com/wuzbak/Unitary-Manifold-}\\[4pt]
\small\itshape Theory, framework, and scientific direction: ThomasCory Walker-Pearson.\\
\small\itshape Code infrastructure, symbolic verification, and document engineering:\\
\small\itshape GitHub Copilot (AI).}

\date{May 2026}
\maketitle

% ============================================================
\begin{abstract}
We present a five-dimensional Kaluza--Klein (KK) inflation model in which a
compact $S^1/\mathbb{Z}_2$ extra dimension topologically constrains the
inflationary winding number to a small discrete set and algebraically fixes the
Chern--Simons level $\kcs$ without free parameters.  Starting from the
$5$D Einstein--Hilbert action with the standard KK block metric (radion $\phi$,
irreversibility 1-form $\Bmu$), we show:
(i)~The $\mathbb{Z}_2$ orbifold restricts the winding number $\nw$ to odd
integers; the Chern--Simons anomaly gap combined with $N_{\mathrm{gen}}=3$ stable
KK matter species constrains $\nw \in \{5,7\}$.
(ii)~The Atiyah--Patodi--Singer (APS) $\eta$-invariant of the boundary Dirac
operator satisfies $\bar\eta(5)=\tfrac{1}{2}$ (non-trivial spin structure) and
$\bar\eta(7)=0$ (trivial), selecting $\nw=5$ on physical grounds
(PHYSICALLY-MOTIVATED; full geometric proof is outstanding work).
(iii)~For the minimum-step braid $(\nw,\nw{+}2)=(5,7)$, the algebraic identity
$\kcs = n_1^2+n_2^2$ (proved) gives $\kcs=74$ with zero free parameters.
(iv)~The braided sound speed $\cs=12/37$ (algebraically derived) suppresses the
tensor-to-scalar ratio to $r_{\mathrm{braided}}\approx 0.0315$, consistent
with BICEP/Keck 2022.
The model simultaneously predicts:
$\ns\approx 0.9635$ (Planck 2018 $< 1\sigma$),
$r\approx 0.0315$ (BICEP/Keck $<0.036$), and
two cosmic birefringence angles $\beta\in\{0.273^\circ, 0.331^\circ\}$
separated by $0.058^\circ = 2.9\,\sigma_{\mathrm{LB}}$.
LiteBIRD ($\sim$2032) provides a decisive discriminating test.
\end{abstract}

% ============================================================
\section{Introduction}
\label{sec:intro}

Inflationary models generically predict a spectral index $\ns$ and a
tensor-to-scalar ratio $r$ consistent with observation, but rarely fix both
simultaneously from topology without continuous free parameters.  Cosmic
birefringence adds a third observable that is sensitive to axion-like couplings.
Fitting all three simultaneously from integer data is a strong constraint.

In this paper we show that a compact $S^1/\mathbb{Z}_2$ Kaluza--Klein dimension
with a braided winding structure:
\begin{itemize}
  \item Restricts $\nw$ to a two-element set $\{5,7\}$ by topology and stability,
  \item Algebraically fixes the Chern--Simons level $\kcs = 5^2+7^2 = 74$,
  \item Predicts $\ns$, $r$, and $\beta$ from these integers alone.
\end{itemize}
The decisive test is LiteBIRD's measurement of $\beta$ to $\pm 0.02^\circ$: both
viable sectors predict distinct birefringence angles separated by
$0.058^\circ\approx 2.9\,\sigma_{\mathrm{LB}}$, so a single observation
discriminates them.

\paragraph{Scope of this paper.}
This paper is narrowly focused on the inflation and CMB predictions of the
Kaluza--Klein sector.  The broader Unitary Manifold framework — including
implications for biology, consciousness, governance, and theology — is
documented in the companion repository
\cite{WalkerPearson2026} but is explicitly excluded here.  Those extensions
are AxiomZero-commissioned implications and not claims of this paper.

\paragraph{Notation.}
We distinguish $\aNM \equiv \phiz^{-2}$ (the nonminimal KK curvature-scalar
coupling from the cross-block Riemann sector $\mathcal{R}^\mu{}_{5\nu5}$) from
$\aem=1/137.036$ (the electromagnetic fine structure constant).
These are physically distinct quantities.

% ============================================================
\section{The 5D Metric Ansatz and Dimensional Reduction}
\label{sec:kk}

\subsection{The 5D action and metric}

We start from the 5D Einstein--Hilbert action:
\begin{equation}
  S_5 = \frac{1}{16\pi G_5}\int d^5x\,\sqrt{-G}\,R_5 ,
  \label{eq:s5}
\end{equation}
with the standard KK block metric (cylinder condition $\partial_5 G_{AB}=0$):
\begin{equation}
  G_{AB} = \begin{pmatrix} g_{\mu\nu} + \lam^2\phi^2 B_\mu B_\nu & \lam\phi B_\mu \\
                            \lam\phi B_\nu & \phi^2 \end{pmatrix} ,
  \label{eq:metric}
\end{equation}
where $g_{\mu\nu}$ is the 4D metric, $\Bmu$ the irreversibility 1-form (KK
vector), $\phi$ the radion (entropic dilaton, $G_{55}=\phi^2$), and $\lam$
the KK coupling.  This is the standard Kaluza--Klein ansatz
\cite{Kaluza1921,Klein1926}; our new content is the physical interpretation
of $\Bmu$ as an irreversibility 1-form and the derivation of the compactification
radius from the fixed-point iteration described in Section~\ref{sec:ftum}.

\subsection{Dimensional reduction}

Integrating $S_5$ over $y\in[0,2\pi R]$ (zero-mode truncation) yields the
4D effective action in the Einstein frame:
\begin{equation}
  S_4 = \frac{1}{16\pi G_4}\int d^4x\,\sqrt{-g}\,\phi
        \Bigl[R_4 - \tfrac{1}{4}\lam^2\phi^2 H_{\mu\nu}H^{\mu\nu}
              + \tfrac{(\partial\phi)^2}{\phi^2}\Bigr]
        + i\int d^4x\,\Bmu J^\mu_{\inf} ,
  \label{eq:s4}
\end{equation}
where $H_{\mu\nu}=\partial_\mu B_\nu - \partial_\nu B_\mu$ and
$J^\mu_{\inf}=\phi^2 u^\mu$ is the information 4-current.

\subsection{The nonminimal coupling $\aNM$}

The cross-block Riemann component $\mathcal{R}^\mu{}_{5\nu5}$ contributes,
after integration, the nonminimal term $\aNM R H_{\mu\nu}H^{\mu\nu}$ with
\begin{equation}
  \aNM = \frac{\lP^2}{L_5^2} = \frac{1}{\phi_0^2} ,
  \label{eq:aNM}
\end{equation}
where $L_5=\phi_0\lP$ is the compactification radius.  This coupling is
\emph{derived}, not inserted; it equals $\phiz^{-2}$ evaluated at the
stabilised radion vev (Section~\ref{sec:ftum}).
\textbf{We emphasise: $\aNM \neq \aem$.}  These are different quantities;
$\aem=1/137.036$ appears only in the Standard Model matching context
\cite{PDG2022}.

\subsection{Recovery of electromagnetism}

The identification $\lam B_\mu \equiv A_\mu$ recovers electromagnetism as the
zero mode of the KK gauge field.  This is the standard KK construction and is
\emph{recovered}, not predicted: it follows from the ansatz~(\ref{eq:metric})
by construction, exactly as in the original Kaluza--Klein programme.  What is
new here is (i) the physical interpretation of $B_\mu$ as the irreversibility
1-form and (ii) the derivation of the compactification scale from the FTUM
rather than as a free parameter.

% ============================================================
\section{Topological Constraint on $\nw$}
\label{sec:topology}

\subsection{$\mathbb{Z}_2$ orbifold: odd winding numbers}

The compact dimension is $S^1/\mathbb{Z}_2$ with the involution $y\to -y$.
KK zero modes must be $\mathbb{Z}_2$-even, which restricts the winding quantum
number to odd integers:
\begin{equation}
  \nw \in \{1,3,5,7,9,\ldots\} .
  \label{eq:z2}
\end{equation}
This is a topological consequence of the orbifold (Pillar~39 of the companion
codebase \cite{WalkerPearson2026}; see also \cite{Horava1996}).

\subsection{CS anomaly gap and $N_{\mathrm{gen}}=3$}

The Chern--Simons anomaly protection gap requires $\Delta_{\mathrm{CS}}=\nw$
stability modes.  Demanding exactly $N_{\mathrm{gen}}=3$ stable KK matter
species gives $\nw\in[4,8]$.  Combined with equation~(\ref{eq:z2}):
\begin{equation}
  \boxed{\nw \in \{5,7\}.}
  \label{eq:nw57}
\end{equation}
This is an algebraic consequence of the orbifold and the generation number;
it is proved in the companion codebase (Pillar~67, function
\texttt{orbifold\_uniqueness()}).

\subsection{APS $\eta$-invariant selection}

The Atiyah--Patodi--Singer $\eta$-invariant of the boundary Dirac operator is
\cite{APS1975}:
\begin{equation}
  \bar\eta(\nw) = \frac{T(\nw)}{2} \bmod 1 ,
  \qquad
  T(\nw) = \frac{\nw(\nw+1)}{2} \quad\text{(triangular number)}.
  \label{eq:eta}
\end{equation}
This formula is derived by three independent analytic methods (Hurwitz
$\zeta$-function, CS inflow, $\mathbb{Z}_2$ zero-mode parity) in
Pillar~70-B of the codebase.  Evaluating:
\begin{equation}
  \bar\eta(5) = \frac{15}{2} \bmod 1 = \tfrac{1}{2}
  \qquad\text{(non-trivial spin structure)},
  \label{eq:eta5}
\end{equation}
\begin{equation}
  \bar\eta(7) = \frac{28}{2} \bmod 1 = 0
  \qquad\text{(trivial spin structure)}.
  \label{eq:eta7}
\end{equation}

\paragraph{Selection of $\nw=5$ (PHYSICALLY-MOTIVATED).}
The Standard Model has left-handed weak-isospin doublets at the orbifold fixed
points.  Left-handed zero modes survive the $\mathbb{Z}_2$ projection only
under $\Omega_{\mathrm{spin}}=-\Gamma^5$, which belongs to the $\bar\eta=\tfrac{1}{2}$
spin-structure class.  Combined with equations~(\ref{eq:eta5})--(\ref{eq:eta7}),
this selects $\nw=5$ uniquely from $\{5,7\}$.
\begin{remark}
  This selection is PHYSICALLY-MOTIVATED: the physics argument is compelling
  and self-consistent, but a purely geometric proof — showing from the 5D metric
  boundary conditions alone that $\Omega_{\mathrm{spin}}=-\Gamma^5$ without
  invoking SM chirality as input — remains outstanding work.  Elevating this to
  PROVED would require that additional step.
\end{remark}

\paragraph{Observational consistency.}
Independently, the Planck 2018 value $\ns=0.9649\pm0.0042$ is consistent with
$\nw=5$ at $0.33\sigma$ and excludes $\nw=7$ at $3.9\sigma$
(computed from $\ns\approx 1-36/\phi_{0,\mathrm{eff}}^2$ with
$\phi_{0,\mathrm{eff}}=\nw\cdot 2\pi\cdot\phi_0$).

% ============================================================
\section{The Braid Algebra: $\kcs$, $\cs$, and Observables}
\label{sec:braid}

\subsection{Algebraic derivation of $\kcs=74$}

\begin{theorem}[Braid CS level, Pillar~58]
  For any braid pair $(n_1,n_2)$ in the $S^1/\mathbb{Z}_2$ KK dimension, the
  effective Chern--Simons level is
  \begin{equation}
    \kcs = n_1^2 + n_2^2 .
    \label{eq:kcs}
  \end{equation}
\end{theorem}
\begin{proof}
  The primary CS level is $k_{\mathrm{prim}}=2(n_1^3+n_2^3)/(n_1+n_2)=
  2(n_1^2-n_1 n_2+n_2^2)$.  The $\mathbb{Z}_2$ correction is
  $\Delta k = (n_2-n_1)^2 = n_1^2-2n_1 n_2+n_2^2$.
  Therefore $\kcs = k_{\mathrm{prim}}-\Delta k = n_1^2+n_2^2$.  $\square$
\end{proof}

For the minimum-step braid $(\nw,\nw{+}2)=(5,7)$:
\begin{equation}
  \boxed{\kcs = 5^2+7^2 = 74 .}
  \label{eq:kcs74}
\end{equation}
This is an algebraic identity; it introduces zero free parameters.
The birefringence observation (Section~\ref{sec:cmb}) confirms $\kcs=74$;
it does not freely tune it.

\subsection{Braided sound speed}

The kinetic mixing parameter is $\rho = 2n_1 n_2/\kcs = 70/74 = 35/37$.
Under the resonance condition~(\ref{eq:kcs}), the canonically-normalised
braided sound speed is:
\begin{equation}
  \cs = \sqrt{1-\rho^2} = \frac{(n_2-n_1)(n_1+n_2)}{\kcs}
      = \frac{2\times 12}{74} = \frac{12}{37} \approx 0.3243 .
  \label{eq:cs}
\end{equation}
This is algebraically derived; it is not a fit.

\subsection{Braided tensor-to-scalar ratio}

In non-canonical kinetic inflation with sound speed $\cs$, the tensor-to-scalar
ratio is suppressed at leading slow-roll order \cite{GarrigaMukhanov1999}:
\begin{equation}
  r_{\mathrm{braided}} = r_{\mathrm{bare}}\times\cs .
  \label{eq:rbraid}
\end{equation}
The bare ratio is $r_{\mathrm{bare}}=96/\phi_{0,\mathrm{eff}}^2\approx 0.097$
(single winding mode, $\nw=5$).  Applying~(\ref{eq:cs}):
\begin{equation}
  r_{\mathrm{braided}} \approx 0.097\times\frac{12}{37} \approx 0.0315 .
  \label{eq:rval}
\end{equation}
\begin{remark}
  Equation~(\ref{eq:rbraid}) is STRONGLY MOTIVATED: it follows Garriga \&
  Mukhanov (1999) for non-canonical kinetic inflation, and $\cs=12/37$ is
  algebraically derived.  A full field-theoretic derivation starting from the
  5D action~(\ref{eq:s5}), canonically normalising the braided kinetic matrix,
  and tracking all slow-roll corrections is outstanding work that would complete
  this argument to DERIVED status.
\end{remark}

% ============================================================
\section{CMB Predictions and Observational Status}
\label{sec:cmb}

\subsection{Three-observable prediction table}

Table~\ref{tab:obs} summarises the predictions.

\begin{table}[h]
\centering
\caption{Predicted CMB observables vs current constraints.}
\label{tab:obs}
\begin{tabular}{lccc}
\toprule
Observable & Prediction & Current data & Status \\
\midrule
$\ns$ & $0.9635$ & $0.9649\pm0.0042$ \cite{Planck2018} & Pass ($0.33\sigma$) \\
$r$ & $0.0315$ & $<0.036$ \cite{BICEPKeck2022} & Pass \\
$\beta$ (sector 5,7) & $0.331^\circ$ & $\sim0.35^\circ$ (hint) \cite{MinamiKomatsu2020,DiegoPalazuelos2022} & Consistent \\
$\beta$ (sector 5,6) & $0.273^\circ$ & $\sim0.35^\circ$ (hint) & Marginal \\
\bottomrule
\end{tabular}
\end{table}

All three observables are consistent with current data.  The birefringence
signal is currently a $2$--$3\,\sigma$ hint from the Planck polarisation data;
it is not a confirmed detection.

\subsection{Cosmic birefringence formula}

The axion-like coupling of $B_\mu$ through the CS term at level $\kcs$ gives:
\begin{equation}
  g_{a\gamma\gamma} = \frac{\kcs\,\aNM}{2\pi^2 r_c} ,
  \label{eq:gcoup}
\end{equation}
where $r_c$ is the compactification radius and $\aNM=\phiz^{-2}$ is the
nonminimal coupling defined in~(\ref{eq:aNM}).
The predicted birefringence angles for the two viable sectors are:
\begin{align}
  \beta_{(5,7)} &\approx 0.331^\circ \quad (\kcs=74) , \\
  \beta_{(5,6)} &\approx 0.273^\circ \quad (\kcs=61=5^2+6^2) .
\end{align}

% ============================================================
\section{The LiteBIRD Falsification Window}
\label{sec:litebird}

The two sectors produce a birefringence gap:
\begin{equation}
  \Delta\beta = 0.331^\circ - 0.273^\circ = 0.058^\circ .
  \label{eq:gap}
\end{equation}
LiteBIRD's projected sensitivity is $\sigma_{\mathrm{LB}}\approx 0.02^\circ$,
so the gap is $0.058^\circ/0.02^\circ \approx 2.9\,\sigma_{\mathrm{LB}}$:
the two sectors are discriminable at $>2\sigma$ by a single LiteBIRD
measurement.

\paragraph{Falsification conditions.}
The braided-winding mechanism is falsified if LiteBIRD measures:
\begin{enumerate}
  \item $\beta$ outside the admissible window $[0.22^\circ, 0.38^\circ]$; or
  \item $\beta$ in the predicted gap $[0.29^\circ, 0.31^\circ]$
        (inconsistent with both sectors simultaneously).
\end{enumerate}
A null detection ($\beta=0$) at LiteBIRD precision would also falsify the
birefringence sector, even if $\ns$ and $r$ remain consistent.

% ============================================================
\section{The Radion Fixed Point and $\phiz$ Self-Consistency}
\label{sec:ftum}

The compactification radius is not a free parameter.  The Fixed-Topology
Unitary Manifold (FTUM) operator $U=I+H+T$ (irreversibility, holographic
projection, time translation) has a unique fixed point $\Psi^*$ satisfying
$U\Psi^*=\Psi^*$, which gives $\phi_0\approx 1$ in Planck units.  Applied
through the KK Jacobian $\mathcal{J}=\nw\cdot 2\pi\sqrt{\phi_0}$:
\begin{equation}
  \phi_{0,\mathrm{eff}} = \nw\times 2\pi\times\phi_0 = 5\times 2\pi\times 1
                        \approx 31.416 .
  \label{eq:phieff}
\end{equation}
The self-consistency condition (Pillar~56, \texttt{phi0\_closure.py}) shows
that three a-priori distinct definitions of $\phi_0$ collapse to the same
value under the $\cs$-corrected slow-roll formula
$\ns=1-36(1+\cs^2)/\phi_{0,\mathrm{eff}}^2$, confirming internal consistency.

% ============================================================
\section{Open Problems}
\label{sec:open}

The following problems are explicitly unresolved and documented here to avoid
overclaiming:

\begin{enumerate}
  \item \textbf{First-principles $c_L$ (Yukawa texture).}
        The Yukawa texture constants $c_L$ determining fermion mass ratios are
        currently derived by bisection at the 5D fixed point (Pillar~98).
        A geometric derivation from 5D orbifold boundary conditions is outstanding.
  \item \textbf{Purely geometric proof of $\nw=5$.}
        The APS argument (Section~\ref{sec:topology}) invokes SM left-handedness.
        A proof from 5D metric boundary conditions alone, without SM input, would
        elevate the selection from PHYSICALLY-MOTIVATED to PROVED.
  \item \textbf{Full derivation of $r_{\mathrm{braided}}=r_{\mathrm{bare}}\times\cs$.}
        Canonical normalisation of the braided kinetic matrix from the 5D action.
  \item \textbf{$\mathrm{SU}(3)\times\mathrm{SU}(2)$ from higher-dimensional compactification.}
        The current 5D framework recovers U(1) electromagnetism by the standard
        KK construction; $\mathrm{SU}(3)\times\mathrm{SU}(2)$ requires a
        minimum of 11D \cite{Witten1981}.
  \item \textbf{Canonical quantisation of $\phi$.}
        The Hamiltonian analysis of the radion sector has not been completed.
\end{enumerate}

% ============================================================
\section{Conclusion}
\label{sec:conclusion}

We have shown that a 5D Kaluza--Klein inflation model with a compact
$S^1/\mathbb{Z}_2$ dimension and a braided winding structure simultaneously
satisfies current CMB constraints on $\ns$, $r$, and (tentatively) $\beta$
through a chain of algebraic and topological arguments with zero continuous
free parameters.

The key results are:
\begin{itemize}
  \item $\nw\in\{5,7\}$ is an algebraic consequence of the orbifold topology
        and the generation number.
  \item $\kcs=74=5^2+7^2$ is an algebraic identity (Theorem~\ref{eq:kcs}),
        not a fit.
  \item $\cs=12/37$ is algebraically derived.
  \item The CMB predictions $\ns\approx 0.9635$, $r\approx 0.0315$, and
        $\beta\in\{0.273^\circ, 0.331^\circ\}$ follow from these integers.
  \item LiteBIRD ($\sim$2032) provides a decisive test: the two viable sectors
        differ in birefringence by $0.058^\circ\approx 2.9\,\sigma_{\mathrm{LB}}$.
\end{itemize}

The framework makes a falsifiable prediction with a clear timeline: if
LiteBIRD measures $\beta$ outside $[0.22^\circ, 0.38^\circ]$ or inside the
gap $[0.29^\circ, 0.31^\circ]$, the braided-winding mechanism is ruled out.
We invite independent verification of all calculations via the public codebase
\cite{WalkerPearson2026}, which provides 15,096 passing tests with 0 failures.

% ============================================================
\section*{Acknowledgements}

Code infrastructure, symbolic verification, and document engineering by
GitHub Copilot (AI).  The author thanks the Planck, BICEP/Keck, and LiteBIRD
collaborations for making their data and specifications public.

% ============================================================
\bibliography{references}

\begin{thebibliography}{99}

\bibitem{Kaluza1921}
T. Kaluza,
\textit{Zum Unit{\"a}tsproblem der Physik},
Sitzungsberichte der Preussischen Akademie der Wissenschaften, 966--972 (1921).

\bibitem{Klein1926}
O. Klein,
\textit{Quantentheorie und f{\"u}nfdimensionale Relativit{\"a}tstheorie},
Zeitschrift f{\"u}r Physik \textbf{37}, 895--906 (1926).

\bibitem{APS1975}
M.F. Atiyah, V.K. Patodi, I.M. Singer,
\textit{Spectral asymmetry and Riemannian geometry I},
Mathematical Proceedings of the Cambridge Philosophical Society
\textbf{77}, 43--69 (1975).

\bibitem{GarrigaMukhanov1999}
J. Garriga, V.F. Mukhanov,
\textit{Perturbations in k-inflation},
Physics Letters B \textbf{458}, 219--225 (1999).
arXiv:hep-th/9904176.

\bibitem{Horava1996}
P. Ho\v{r}ava, E. Witten,
\textit{Heterotic and Type I string dynamics from eleven dimensions},
Nuclear Physics B \textbf{460}, 506--524 (1996).

\bibitem{Witten1981}
E. Witten,
\textit{Search for a realistic Kaluza--Klein theory},
Nuclear Physics B \textbf{186}, 412--428 (1981).

\bibitem{Planck2018}
Planck Collaboration (N. Aghanim et al.),
\textit{Planck 2018 results VI: Cosmological parameters},
Astronomy \& Astrophysics \textbf{641}, A6 (2020).
arXiv:1807.06209.

\bibitem{BICEPKeck2022}
BICEP/Keck Collaboration,
\textit{BicepKeck XVIII: Improved constraints on primordial gravitational waves
using Keck Array and BICEP3 through the 2021 observing season},
Physical Review Letters (2023). arXiv:2212.05560.

\bibitem{MinamiKomatsu2020}
Y. Minami, E. Komatsu,
\textit{New extraction of the cosmic birefringence from the Planck 2018
polarization data},
Physical Review Letters \textbf{125}, 221301 (2020).
arXiv:2011.11461.

\bibitem{DiegoPalazuelos2022}
J. Diego-Palazuelos et al.,
\textit{Cosmic birefringence from the Planck data},
Physical Review Letters \textbf{128}, 091302 (2022).
arXiv:2201.07682.

\bibitem{PDG2022}
Particle Data Group (R.L. Workman et al.),
\textit{Review of Particle Physics},
Progress of Theoretical and Experimental Physics \textbf{2022}, 083C01 (2022).

\bibitem{WalkerPearson2026}
T. Walker-Pearson,
\textit{The Unitary Manifold: A 5D Gauge Geometry of Emergent Irreversibility},
v9.29, Zenodo, \url{https://doi.org/10.5281/zenodo.19584531} (2026).
Code: \url{https://github.com/wuzbak/Unitary-Manifold-}.

\end{thebibliography}

\end{document}
